\subsubsection{Пользовательский интерфейс}

Веб-приложение используется для запуска эксперимента. Пользователь выбирает
модель хранения, профиль нагрузки, \texttt{seed}, длительность ступени, число
параллельных клиентов и размеры пакетов. После этого приложение выполняет
очистку, подготовку данных и нагрузочный прогон.

Все команды проходят через Express API. Это позволяет запускать только один
прогон за раз и не смешивать результаты разных моделей.

\subsubsection{Отображение метрик}

Для просмотра результатов в интерфейс встроена панель Grafana
\cite{grafana_database_observability_monitor}. В отчете используются
следующие показатели: общее количество операций, доля ошибок, пропускная
способность, объем обработанных данных, p95 и p99 задержек, пропускная
способность по моделям и ступеням нагрузки, а также задержки операций записи
и чтения. Эти метрики СУБД используются для построения вывода о производительности
конкретной моделей хранения данных.
Метрики сравниваются между собой на разных моделях на разных уровнях нагрузки и
сценариях.
